\documentclass[11pt]{article}

\usepackage[T1]{fontenc}
\usepackage{lmodern}
\usepackage{amsmath,amssymb,amsthm}
\usepackage[margin=1in]{geometry}
\usepackage{booktabs}
\usepackage{enumitem}
\usepackage{microtype}
\usepackage[hidelinks]{hyperref}

\hypersetup{
  pdftitle={General-m minimal-equilibrium stability: mathematical proofs and an exact subcritical steady obstruction}
}

\newcommand{\R}{\mathbb{R}}
\newcommand{\norm}[1]{\left\lVert #1\right\rVert}
\newcommand{\abs}[1]{\left\lvert #1\right\rvert}
\setlength{\emergencystretch}{3em}

\newtheorem{theorem}{Theorem}
\newtheorem{proposition}{Proposition}
\newtheorem{lemma}{Lemma}
\newtheorem{remark}{Remark}

\title{General-\(m\) Minimal-Equilibrium Stability\\
\large Explicit sufficient regions, the bifurcation boundary, and an exact
\(m=3\) counterexample}
\author{}
\date{27 July 2026}

\begin{document}
\maketitle

\begin{abstract}
We study the one-dimensional Neumann minimal chemotaxis system with mobility
\(u^m\).  The sufficient stability conditions are written explicitly.  For
\(1\le m<2\), two energy routes give global attraction on the conserved-mass
hyperplane.  At \(m=2\), the same conclusion holds under a sharp
parameter-level endpoint admissibility condition.  For every \(m>1\), a
negative-power oscillation estimate gives a separate small-sensitivity
theorem.

The boundary \(m=2\) belongs to the energy method, not to the dynamics.  A
weakly nonlinear calculation determines the direction of the first Neumann
steady branch.  For
\[
 m=3,\qquad \beta=\gamma=\mu=\nu=U=1,
\]
a Banach-space implicit-function argument constructs a positive nonconstant
steady state at a sensitivity strictly below the linear threshold.  This
steady state is an exact obstruction to global attraction.  The proof is
analytic and does not rely on numerical continuation.
\end{abstract}

\tableofcontents

\section{Model and meaning of stability}

On \(I=(0,1)\), consider
\begin{align}
 u_t
 &=u_{xx}-\chi_0\partial_x
   \left(\frac{u^m}{(1+v)^\beta}v_x\right),             \label{eq:u}\\
 0&=v_{xx}-\mu v+\nu u^\gamma,                          \label{eq:v}\\
 u_x(t,0)&=u_x(t,1)=v_x(t,0)=v_x(t,1)=0.                \label{eq:bc}
\end{align}
All parameters satisfy
\[
 m>0,\quad \beta\ge1,\quad \gamma,\mu,\nu>0,\quad
 \chi_0>0.
\]
Integration of \eqref{eq:u} gives the conserved mass
\[
 \int_0^1u(t,x)\,dx=U>0.
\]
The constant equilibrium on this mass hyperplane is
\begin{equation}
 u_*=U,\qquad v_*=\frac{\nu}{\mu}U^\gamma.              \label{eq:eq}
\end{equation}

Here ``global attraction'' means the following statement about every
positive, bounded, global classical solution with mass \(U\):
\[
 \norm{u(t)-U}_{L^\infty}\longrightarrow0,
\]
and there are orbit-dependent constants \(C,\rho,t_0>0\) such that
\begin{equation}
 \norm{u(t)-U}_{C^1}+\norm{v(t)-v_*}_{C^1}
 \le Ce^{-\rho t}\qquad(t\ge t_0).                     \label{eq:eventualC1}
\end{equation}
The positive entry time is necessary when the initial datum is only
continuous.

\section{Linear stability threshold}

Let \(\lambda_n=n^2\pi^2\) be the Neumann eigenvalues.  Write
\[
 u=U+p,\qquad v=v_*+q.
\]
For the \(n\)-th nonconstant mode, \eqref{eq:v} gives
\[
 \widehat q_n
 =\frac{\nu\gamma U^{\gamma-1}}{\mu+\lambda_n}
   \widehat p_n.
\]
Linearization of \eqref{eq:u} therefore gives the growth rate
\begin{equation}
 \sigma_n(\chi_0)
 =-\lambda_n+
 \chi_0\frac{\nu\gamma U^{m+\gamma-1}\lambda_n}
 {(1+v_*)^\beta(\mu+\lambda_n)}.                       \label{eq:sigma}
\end{equation}

\begin{proposition}[linear threshold]
The constant equilibrium is linearly stable on the mass-zero subspace if
\begin{equation}
 \chi_0<\chi_{\rm lin}:=
 \frac{(1+v_*)^\beta(\mu+\pi^2)}
 {\nu\gamma U^{m+\gamma-1}}.                           \label{eq:chilin}
\end{equation}
At equality the first cosine mode is neutral, and it is the first
nonconstant mode to lose stability.
\end{proposition}

\begin{proof}
For \(n\ge1\), \(\sigma_n<0\) is equivalent to
\[
 \chi_0<
 \frac{(1+v_*)^\beta(\mu+\lambda_n)}
 {\nu\gamma U^{m+\gamma-1}}.
\]
The right-hand side is strictly increasing with \(\lambda_n\), so its
minimum occurs at \(n=1\).
\end{proof}

For the normalized tuple \(U=\beta=\gamma=\mu=\nu=1\),
\begin{equation}
 \chi_{\rm lin}=2(1+\pi^2)\approx21.73920880.           \label{eq:fixedlin}
\end{equation}

\section{Explicit sufficient stability conditions}

\subsection{Constants appearing in the two energy routes}

Define the sharp signal saturation factor
\begin{equation}
 \sigma_\beta=
 \sup_{z\ge0}\frac{z}{(1+z)^\beta}
 =
 \begin{cases}
 1,&\beta=1,\\[1mm]
 \dfrac{(\beta-1)^{\beta-1}}{\beta^\beta},&\beta>1,
 \end{cases}                                           \label{eq:sigmabeta}
\end{equation}
and
\begin{equation}
 \chi_\beta=\frac{2(2\beta-1)}{\max\{2,\gamma\}}.       \label{eq:chibeta}
\end{equation}
The formula for \(\sigma_\beta\) follows by differentiating
\(z(1+z)^{-\beta}\); for \(\beta>1\) its unique maximum is
\(z=(\beta-1)^{-1}\).

The energy estimates use an eventual box.  Namely, choose
\(\overline u>0\) and \(\underline v>0\), depending only on the parameters
and \(U\), such that every orbit under consideration satisfies
\begin{equation}
 u(t,x)\le\overline u,\qquad v(t,x)\ge\underline v
 \quad\text{for all sufficiently large \(t\), uniformly in \(x\)}.
                                                               \label{eq:box}
\end{equation}
For \(1\le m<2\), standard one-dimensional \(L^P\) recurrence, Agmon
interpolation, the conserved mass, and positivity of the Neumann elliptic
Green function produce such orbit-independent constants.  The same
construction works at \(m=2\) under the endpoint condition
\begin{equation}
 \chi_0\sqrt\mu\,\sigma_\beta U\max\{2,\gamma\}<1.      \label{eq:critical}
\end{equation}
The role of \eqref{eq:critical} is to leave room for a restarted exponent
strictly larger than \(2\) and at least \(\gamma\).

On the interval \([0,\overline u]\), define the power slope
\begin{equation}
 L_\gamma(U,\overline u)=
 \begin{cases}
 U^{\gamma-1},&0<\gamma\le1,\\[1mm]
 \gamma\overline u^{\gamma-1},&\gamma>1.
 \end{cases}                                           \label{eq:slope}
\end{equation}
Then
\[
 \abs{s^\gamma-U^\gamma}
 \le L_\gamma(U,\overline u)\abs{s-U},
 \qquad 0\le s\le\overline u.
\]

The entropy threshold is
\begin{equation}
 \chi_E=
 \frac{2\sqrt\mu\,(1+\underline v)^\beta\overline u^{-m}}
 {\nu L_\gamma(U,\overline u)},                        \label{eq:chiE}
\end{equation}
and the signal-energy threshold is
\begin{equation}
 \chi_S=
 \min\left\{
 \frac{\chi_\beta}{2},\
 \sqrt{\chi_\beta},\
 \frac{\mu(1+\underline v)^\beta}
 {\nu\overline u^m}
 \right\}.                                             \label{eq:chiS}
\end{equation}

\begin{theorem}[two explicit stability branches]
Assume \(1\le m<2\), and let \(\overline u,\underline v\) satisfy
\eqref{eq:box}.  The equilibrium \eqref{eq:eq} is globally attracting if
either
\begin{equation*}
 0<\chi_0<\min\{\chi_{\rm lin},\chi_E\},                \tag{E}
\end{equation*}
or
\begin{equation*}
 \gamma=1,\qquad
 0<\chi_0<\min\{\chi_{\rm lin},\chi_S\}.                \tag{S}
\end{equation*}
At \(m=2\), the same conclusion holds under
\eqref{eq:critical} and either \emph{(E)} or \emph{(S)}.
\end{theorem}

\subsection{Proof of the entropy branch}

Define
\begin{equation}
 H_m(s)=\int_U^s
 \left[1-\left(\frac{U}{\tau}\right)^{2m-1}\right]d\tau,
 \qquad
 \mathcal E_u(t)=\int_0^1H_m(u(t,x))\,dx.               \label{eq:entropy}
\end{equation}
For \(m\ge1\), \(H_m\ge0\) and \(H_m(s)=0\) only at \(s=U\).
Put
\[
 c_m=(2m-1)U^{2m-1},
\quad
 G=\int_0^1u^{-2m}u_x^2,
\quad
 X=\int_0^1\frac{u^{-m}u_xv_x}{(1+v)^\beta}.
\]
Differentiation of \eqref{eq:entropy}, followed by integration by parts,
gives
\begin{equation}
 \mathcal E_u'(t)\le-c_mG+\chi_0c_mX.                  \label{eq:Eslope}
\end{equation}
The boundary terms vanish because of \eqref{eq:bc}.

Young's inequality, with
\[
 A=u^{-m}u_x,\qquad B=(1+v)^{-\beta}v_x,
\]
gives
\begin{equation}
 -c_mG+\chi_0c_mX
 \le-\frac{c_m}{2}G+
 \frac{\chi_0^2c_m}{2}
 \int_0^1\frac{v_x^2}{(1+v)^{2\beta}}.                 \label{eq:halfyoung}
\end{equation}
Inside the eventual box, mass conservation and the Neumann Poincaré
inequality imply
\begin{equation}
 G\ge\overline u^{-2m}
 \int_0^1(u-U)^2.                                      \label{eq:Gcoercive}
\end{equation}
The elliptic equation, tested against \(v-v_*\), and the sharp
one-dimensional multiplier bound give
\begin{equation}
 \int_0^1\frac{v_x^2}{(1+v)^{2\beta}}
 \le
 \frac{\nu^2}
 {4\mu(1+\underline v)^{2\beta}}
 \int_0^1(u^\gamma-U^\gamma)^2.                        \label{eq:ellbound}
\end{equation}
Using \eqref{eq:slope},
\[
 \int_0^1(u^\gamma-U^\gamma)^2
 \le L_\gamma(U,\overline u)^2\int_0^1(u-U)^2.
\]
Consequently,
\begin{equation}
 \mathcal E_u'(t)
 \le-C_E\int_0^1(u-U)^2,                               \label{eq:Ediss}
\end{equation}
where
\begin{equation}
 C_E=\frac{c_m}{2}\left[
 \overline u^{-2m}
 -\frac{\chi_0^2\nu^2L_\gamma(U,\overline u)^2}
 {4\mu(1+\underline v)^{2\beta}}
 \right].                                              \label{eq:CE}
\end{equation}
Condition \(\chi_0<\chi_E\) is exactly \(C_E>0\).

Since \(\mathcal E_u\ge0\), integration of \eqref{eq:Ediss} gives
\[
 \int_T^\infty\norm{u(t)-U}_{L^2}^2\,dt<\infty.
\]
Thus there are arbitrarily late times \(t_j\) with
\(\norm{u(t_j)-U}_{L^2}\to0\).  Uniform late-time parabolic H\"older bounds
upgrade this to
\begin{equation}
 \norm{u(t_j)-U}_{L^\infty}\longrightarrow0.           \label{eq:basinentry}
\end{equation}
Indeed, otherwise a point with deviation at least \(\varepsilon\), together
with a uniform spatial modulus of continuity, would produce an interval of
fixed positive length on which the deviation is at least
\(\varepsilon/2\), contradicting the \(L^2\) convergence.

Finally, \(\chi_0<\chi_{\rm lin}\) supplies a spectral gap on the mass-zero
subspace.  The nonlinear remainder is quadratic in a fractional domain
embedded in \(C^1\).  The analytic semigroup estimate and variation of
constants give, for sufficiently small data,
\[
 \norm{(u-U,v-v_*)(t)}_{C^1}
 \le Ce^{-\rho(t-t_j)}
 \norm{u(t_j)-U}_{L^\infty}.
\]
Parabolic smoothing during a fixed positive time connects the
\(L^\infty\)-small slice \eqref{eq:basinentry} to that fractional domain.
This proves \eqref{eq:eventualC1}; the supremum convergence follows
immediately.

\subsection{Proof of the signal-energy branch}

Assume \(\gamma=1\).  Then \(v_*=\nu U/\mu\), and integration of the elliptic
equation gives
\[
 \int_0^1(v-v_*)\,dx=0.
\]
Define
\begin{equation}
 \mathcal E_v(t)=\int_0^1
 \left[v_x^2+\mu(v-v_*)^2\right]dx.                    \label{eq:signalE}
\end{equation}
Differentiating the elliptic equation in time, pairing with \(v-v_*\), and
using the population equation gives the exact identity
\begin{align}
 \mathcal E_v'(t)
 ={}&-2\int_0^1v_{xx}^2
     -2\mu\int_0^1v_x^2 \notag\\
 &\quad
 +2\nu\chi_0\int_0^1
 \frac{u^m v_x^2}{(1+v)^\beta}.                        \label{eq:signalid}
\end{align}
On the eventual box,
\[
 \mathcal E_v'(t)
 \le-2\int_0^1v_{xx}^2-2C_S\int_0^1v_x^2,
\quad
 C_S=\mu-\frac{\nu\chi_0\overline u^m}
 {(1+\underline v)^\beta}.                             \label{eq:CS}
\]
The last entry of \eqref{eq:chiS} makes \(C_S>0\).  Since the signal
perturbation has zero mean, Poincaré's inequality gives
\[
 \mathcal E_v(t)\le(\mu+1)\int_0^1v_x^2.
\]
Hence \(\mathcal E_v\) decays exponentially.  Integration of
\eqref{eq:signalid} also supplies arbitrarily late slices with small
\(\norm{v_{xx}}_{L^2}\).  Because
\[
 u-U=\frac{\mu(v-v_*)-v_{xx}}{\nu},
\]
such slices are \(L^2\)-close to \(U\); smoothing and the one-dimensional
embedding give the same \(L^\infty\) basin entry as in
\eqref{eq:basinentry}.  The linear spectral bootstrap then proves
\eqref{eq:eventualC1}.  The first two entries in \eqref{eq:chiS} are the
smallness bounds used to produce the eventual box and its persistence
estimates.

\subsection{Why the same proof reaches \(m=2\)}

The entropy and signal-energy arguments above use only \(m\ge1\) after the
eventual box has been obtained.  The restriction \(m<2\) enters earlier, in
the \(L^P\) recurrence that produces \eqref{eq:box}.  At \(m=2\), the
recurrence is critical.  Condition \eqref{eq:critical} makes the coefficient
of the restarted moment positive for an exponent \(P>\max\{2,\gamma\}\).
The endpoint Agmon estimate then yields \eqref{eq:box}, after which the same
calculations \eqref{eq:Eslope}--\eqref{eq:CS} apply without change.  Thus
\(m=2\) is an endpoint of this a priori estimate, not a dynamical
bifurcation value.

\section{Every \(m>1\): a separate small-sensitivity route}

For \(m>1\), define
\begin{equation}
 G(t)=\int_0^1u(t,x)^{-2m}u_x(t,x)^2\,dx.               \label{eq:negG}
\end{equation}
The negative-power entropy identity and the sharp saturation inequality
\[
 \frac{v}{(1+v)^\beta}\le\sigma_\beta
\]
give a parameter-only differential inequality.  Indeed, the elliptic
log-gradient bound
\[
 \abs{v_x}\le\sqrt\mu\,v
\]
implies
\[
 \int_0^1\frac{v_x^2}{(1+v)^{2\beta}}
 \le\mu\sigma_\beta^2.
\]
The half-Young estimate \eqref{eq:halfyoung} therefore yields
\begin{equation}
 \mathcal E_u'(t)
 \le-\frac{c_m}{2}
 \left(G(t)-\chi_0^2\mu\sigma_\beta^2\right).           \label{eq:goodslope}
\end{equation}
It follows that, for every \(T>0\) and
\(\varepsilon>0\), there is \(t\ge T\) such that
\begin{equation}
 G(t)<\chi_0^2\mu\sigma_\beta^2+\varepsilon.            \label{eq:goodslice}
\end{equation}
Otherwise \eqref{eq:goodslope} would force \(\mathcal E_u\) to decrease at a
fixed negative rate for all \(t\ge T\), contradicting
\(\mathcal E_u(t)\ge0\).

\begin{lemma}[negative-power oscillation]
For every positive \(C^1\) profile \(u\) on \([0,1]\),
\begin{equation}
 \abs{u(x)^{1-m}-u(z)^{1-m}}
 \le(m-1)\sqrt{G}
 \qquad(x,z\in[0,1]).                                  \label{eq:osc}
\end{equation}
\end{lemma}

\begin{proof}
The fundamental theorem of calculus and Cauchy--Schwarz give
\begin{align*}
 \abs{u(x)^{1-m}-u(z)^{1-m}}
 &\le(m-1)\int_0^1u^{-m}\abs{u_x}\\
 &\le(m-1)
 \left(\int_0^1u^{-2m}u_x^2\right)^{1/2}.
\end{align*}
\end{proof}

For \(0<\delta<U\), set
\begin{equation}
 r_m(U,\delta)=
 \min\left\{
 U^{1-m}-(U+\delta)^{1-m},\
 \frac{(U-\delta)^{1-m}-U^{1-m}}2
 \right\}.                                             \label{eq:radius}
\end{equation}
Both entries are positive because \(1-m<0\).  Continuity and the mass
constraint provide a point \(z\) with \(u(z)=U\).  Hence
\eqref{eq:osc} implies
\begin{equation}
 (m-1)\sqrt G<r_m(U,\delta)
 \quad\Longrightarrow\quad
 \norm{u-U}_{L^\infty}<\delta.                         \label{eq:powerbox}
\end{equation}
To verify \eqref{eq:powerbox}, note that \(s\mapsto s^{1-m}\) is strictly
decreasing.  A point with \(u\ge U+\delta\) or \(u\le U-\delta\) would force
the oscillation from \(U^{1-m}\) to be at least one of the two quantities
used in \eqref{eq:radius}.

\begin{theorem}[all-\(m>1\) small-sensitivity branch]
Assume \(m>1\), \(\chi_0<\chi_{\rm lin}\), and let
\(\delta\in(0,U)\) be any \(L^\infty\) radius contained in the local
spectral basin of the equilibrium.  If
\begin{equation}
 (m-1)\chi_0\sqrt\mu\,\sigma_\beta<r_m(U,\delta),       \label{eq:small}
\end{equation}
then the equilibrium is globally attracting among positive bounded
mass-\(U\) solutions.
\end{theorem}

\begin{proof}
Choose \(\varepsilon>0\) so small that
\[
 (m-1)\sqrt{\chi_0^2\mu\sigma_\beta^2+\varepsilon}
 <r_m(U,\delta).
\]
At the good time supplied by \eqref{eq:goodslice}, the oscillation estimate
and \eqref{eq:powerbox} put the solution inside the local
\(L^\infty\)-basin.  Restarting the autonomous equation at that time and
using the linearly stable semigroup bootstrap proves
\eqref{eq:eventualC1}.
\end{proof}

Equivalently, the explicit threshold associated with \(\delta\) is
\[
 \chi_0<
 \frac{r_m(U,\delta)}
 {(m-1)\sqrt\mu\,\sigma_\beta}.
\]
Linear stability guarantees the existence of at least one positive local
basin radius \(\delta<U\).

\section{Steady-state bifurcation}

\subsection{Integrated steady equation}

Specialize to
\[
 \beta=\gamma=\mu=\nu=1.
\]
At a positive steady state, integration of the zero-flux equation gives
\[
 u_x-\chi\frac{u^m}{1+v}v_x=0.
\]
If \(H_m'(u)=u^{-m}\), then
\[
 H_m(u)=
 \begin{cases}
 \log u,&m=1,\\[1mm]
 \dfrac{u^{1-m}}{1-m},&m\ne1,
 \end{cases}
\]
and every positive steady state satisfies
\begin{equation}
 H_m(u)-\chi\log(1+v)=C,\qquad
 -v_{xx}+v=u,\qquad \int_0^1u=U.                       \label{eq:integrated}
\end{equation}

\subsection{Weakly nonlinear direction coefficient}

Let
\[
 e_1(x)=\cos(\pi x),\qquad e_2(x)=\cos(2\pi x),
\]
and set
\[
 r_1=\frac1{1+\pi^2},\qquad
 r_2=\frac1{1+4\pi^2},\qquad b=1+U.
\]
Seek a mass-preserving branch
\begin{align}
 u&=U+\varepsilon e_1+\varepsilon^2a_2e_2+O(\varepsilon^3), \label{eq:uexp}\\
 v&=U+\varepsilon r_1e_1+\varepsilon^2a_2r_2e_2
       +O(\varepsilon^3),                              \label{eq:vexp}\\
 \chi&=\chi_{\rm lin}+\varepsilon^2\chi_2
       +O(\varepsilon^3).                              \label{eq:chiexp}
\end{align}
Substitution into \eqref{eq:integrated} gives, successively,
\begin{align}
 \chi_{\rm lin}&=\frac{b}{r_1U^m},                     \label{eq:wlin}\\
 a_2&=\frac{m/U-r_1/b}{4(1-r_2/r_1)},                  \label{eq:a2}\\
 \chi_2&=\chi_{\rm lin}S(m,U),                         \label{eq:chi2}\\
 S(m,U)
 &=\frac{a_2}{2}\left(\frac{r_2}{b}-\frac mU\right)
   +\frac{m(m+1)}{8U^2}
   -\frac{r_1^2}{4b^2}.                               \label{eq:S}
\end{align}
The order-\(\varepsilon\) equation yields \eqref{eq:wlin}; the
\(e_2\)-projection at order \(\varepsilon^2\) yields \eqref{eq:a2}; and the
\(e_1\)-solvability condition at order \(\varepsilon^3\) yields
\eqref{eq:chi2}--\eqref{eq:S}.

Thus \(S>0\) gives a forward branch and \(S<0\) gives a backward branch.
The roots of the explicit function \(S\) are approximately
\[
\begin{array}{c|c}
U&m_c(U)\\ \hline
0.5&2.866754572143\\
1&2.936658202934\\
2&3.005668668328
\end{array}
\]
within the normalized family.  These values describe only the direction of
the local steady branch; they do not constitute a global stability diagram.

\section{An exact subcritical steady state at \(m=3\)}

\subsection{Banach-space branch construction}

Fix
\[
 m=3,\qquad \beta=\gamma=\mu=\nu=U=1.
\]
Let \(\mathcal A_2\) be the real even Wiener space
\[
 f(x)=\sum_{n\in\mathbb Z}f_ne^{in\pi x},\qquad
 f_{-n}=f_n\in\R,\qquad
 \norm f_{\mathcal A_2}
 =\sum_{n\in\mathbb Z}(1+\abs n)^2\abs{f_n}<\infty.
\]
It is a Banach algebra and embeds continuously into \(C^2\).
Let \(R=(-\partial_{xx}+1)^{-1}\), so
\[
 R e_n=\frac1{1+n^2\pi^2}e_n.
\]
Restrict to the mean-zero subspace and write
\[
 h=e_1+z,\qquad z\perp e_1.
\]
For amplitude \(a\), set
\[
 u=1+ah,\qquad v=1+aRh.
\]
After dividing the steady flux equation by \(a\), define
\begin{equation}
 \mathcal F(a,z,\chi)=
 h-\chi\,\partial_x^{-1}
 \left[\frac{(1+ah)^3}{2+aRh}\,\partial_xRh\right].
                                                               \label{eq:F}
\end{equation}
The inverse derivative is taken on odd mean-zero Fourier series.  Near
\((a,z,\chi)=(0,0,2(1+\pi^2))\), the denominator is invertible in the
Banach algebra, so \(\mathcal F\) is \(C^3\).

At the base point, the derivative in the complement direction is
\[
 L=I-(1+\pi^2)R.
\]
On the \(n\)-th cosine mode,
\begin{equation}
 Le_n=
 \frac{(n^2-1)\pi^2}{1+n^2\pi^2}e_n.                  \label{eq:Lmode}
\end{equation}
Its kernel on the mean-zero space is precisely the first cosine mode, which
has been removed from \(z\).  For \(n\ge2\), the reciprocal multiplier is
\[
 \frac{1+n^2\pi^2}{(n^2-1)\pi^2},
\]
uniformly bounded in \(n\), so \(L\) has a bounded inverse on the complement.
The derivative in the sensitivity direction is
\[
 -\frac1{2(1+\pi^2)}e_1,
\]
which spans the missing first-mode line.  Hence
\[
 D_{(z,\chi)}\mathcal F
 \colon \{e_1\}^{\perp}\times\R
 \longrightarrow \{\text{mean-zero even profiles}\}
\]
is a bounded isomorphism.

The Banach implicit-function theorem now gives \(a_0>0\) and a \(C^3\)
branch
\[
 a\longmapsto\bigl(z(a),\chi(a)\bigr),
 \qquad
 z(0)=0,\quad \chi(0)=2(1+\pi^2),
\]
for \(\abs a<a_0\), satisfying \(\mathcal F=0\).

\subsection{Exact branch direction}

Multiplying \eqref{eq:F} by the signal denominator and differentiating gives
the polynomial flux identity
\begin{equation}
 (2+aRh)\,\partial_xh
 =\chi(a)(1+ah)^3\partial_xRh.                          \label{eq:poly}
\end{equation}
Projecting the first derivative of \eqref{eq:poly} onto \(e_1\) gives
\(\chi'(0)=0\).  The mode \(e_2\) determines \(z'(0)\), equivalently the
coefficient \(a_2\) in \eqref{eq:a2}.  Projecting the second derivative onto
\(e_1\), or substituting \(m=3,U=1\) into
\eqref{eq:a2}--\eqref{eq:S}, gives
\begin{equation}
 \chi''(0)=
 -\frac{6\pi^4+37\pi^2+25}
 {24\pi^2(\pi^2+1)}<0.                                \label{eq:second}
\end{equation}
Every factor in the denominator is positive and the numerator is positive,
so the sign is strict.

By continuity of \(\chi''\), after reducing \(a_0\) if necessary,
\(\chi'(a)<0\) for \(0<a<a_0\).  Since \(\chi(0)>0\), there is
\(a_1\in(0,a_0)\) such that
\begin{equation}
 0<\chi(a)<2(1+\pi^2)
 \qquad(0<a<a_1).                                      \label{eq:backward}
\end{equation}

\subsection{Positivity and the dynamical obstruction}

The embedding \(\mathcal A_2\hookrightarrow C^2\) and
\(z(a)\to0\) show that, for small \(a>0\),
\[
 u_a=1+a(e_1+z(a))>0,\qquad
 v_a=1+aR(e_1+z(a))>0.
\]
The cosine representation gives the Neumann boundary conditions, the
missing zero Fourier coefficient gives
\[
 \int_0^1u_a\,dx=1,
\]
and the normalized first coefficient makes \(u_a\) nonconstant.
Differentiating \eqref{eq:poly} recovers
\begin{align*}
 0&=(u_a)_{xx}
 -\chi(a)\partial_x\left(\frac{u_a^3(v_a)_x}{1+v_a}\right),\\
 0&=(v_a)_{xx}-v_a+u_a.
\end{align*}

\begin{theorem}[subcritical nonconstant steady state]
There exist \(0<\chi<2(1+\pi^2)\) and positive, nonconstant \(C^2\)
Neumann profiles \(u,v\), of mass one, satisfying the \(m=3\) steady system.
\end{theorem}

Regard these profiles as constant-in-time functions.  They are a positive,
bounded, global classical orbit.  If the constant equilibrium \((1,1)\)
attracted every such orbit, then the time-independent function \(u\) would
converge uniformly to \(1\), forcing \(u\equiv1\), contrary to
nonconstancy.

\begin{theorem}[failure of global attraction below the linear threshold]
For \(m=3\), \(\beta=\gamma=\mu=\nu=U=1\), there is a sensitivity
\(0<\chi<\chi_{\rm lin}\) for which the constant equilibrium is linearly
stable but not globally attracting on the mass-one hyperplane.
\end{theorem}

\section{Resolved region and remaining gap}

The proved picture is:
\begin{enumerate}[leftmargin=2em]
\item for \(1\le m<2\), the explicit entropy or signal-energy conditions
  \emph{(E)}--\emph{(S)} imply global attraction;
\item \(m=2\) is included under \eqref{eq:critical};
\item for every \(m>1\), \eqref{eq:small} gives a separate
  small-sensitivity stability region;
\item at \(m=3,U=1\), a nonconstant positive steady state exists strictly
  below the linear threshold, so linear stability does not imply global
  attraction.
\end{enumerate}

No necessary-and-sufficient global stability diagram is claimed in the
region between the sufficient thresholds and the backward steady branch.
In particular, \(m<m_c(U)\) only means that the first local branch points
forward; it does not prove global attraction all the way to
\(\chi_{\rm lin}\).  Conversely, \(m>m_c(U)\) predicts a local backward
branch, but the analytic branch construction above was carried out for the
fixed \(m=3,U=1\) tuple.

\end{document}
